\casechapter{Combinatorics as Dimension}{From Oddtown to cap sets}{Linearization}{ch:dimension}
\begin{refsection}[chapters/09-dimension.bib]

\newcommand{\dimensionSR}{\operatorname{sr}}

\begin{chaptersummary}
How many clubs can a town have if every club has an odd number of members and every two clubs share an even number? How small can a set be that contains a line in every direction? How many cards of the game SET can lie on a table with no SET among them? For decades such questions were fought with counting and geometry, which gave weak bounds and, for Borsuk's partition problem, a sixty-year-old conjecture that turned out to be false. The move is to linearize: send each object to a vector or a polynomial in a space of known dimension, so that the defining property becomes linear independence or forces a low-degree polynomial to vanish where it cannot. Then Fisher's inequality takes two sentences, Borsuk's conjecture fails in dimension $1325$, the finite-field Kakeya problem is settled in half a page, and cap sets are exponentially small.
\end{chaptersummary}

%==========================================================================
\section{Before: counting by hand}\label{sec:dimension:before}

Each of the five problems below asks how large, or how small, a finite configuration can be when it obeys a local rule about how two sets meet, which directions a set contains, or which triples of points lie on a line. Each was first attacked with the tools of its own field, which found the truth in some cases, gave weak bounds in others, and in one case supported a false conjecture.

\subsection{Clubs in Oddtown}\label{sec:dimension:oddtown-before}

\begin{problem}[Oddtown]\label[problem]{prob:dimension:oddtown}
A town has $n$ citizens, who form clubs. Every club has an odd number of members, and every two distinct clubs have an even number of members in common. How many clubs can there be?
\end{problem}

The $n$ one-member clubs obey the rules. Is $n$ the maximum? The difficulty shows in a neighboring town. In Eventown every club, and every intersection of two clubs, has an even number of members. Pair off the citizens into $\floor{n/2}$ couples and let every union of couples, the empty one included, be a club: the rules hold, and there are $2^{\floor{n/2}}$ clubs. One parity in the rules separates a town with $n$ clubs from a town with exponentially many, so a proof for Oddtown must use the oddness of the clubs, and counting arguments are good at sizes and bad at parities.\index{Oddtown theorem} Both bounds appear in the answer, by \textcite{Berlekamp1969} and independently by Graver, to a question of Erd\H{o}s about families with even pairwise intersections.

\subsection{Fisher's inequality, by counting}\label{sec:dimension:fisher-before}

In 1940 R.~A. Fisher studied balanced incomplete block designs, a tool of experimental design: $v\ge2$ varieties are tested in $b$ blocks, each block contains $k<v$ varieties, and every two varieties appear together in exactly $\lambda\ge1$ blocks.\index{Fisher's inequality}

\begin{theorem}[\textcite{Fisher1940}]\label{thm:dimension:fisherbibd}
Every balanced incomplete block design has at least as many blocks as varieties: $b\ge v$.
\end{theorem}

\begin{proof}
The plan, which is Fisher's, is to fix one block and apply the Cauchy--Schwarz inequality to its intersections with the other blocks. Counting the pairs formed by another variety and a block containing both shows that every variety lies in the same number $r$ of blocks, with $r(k-1)=\lambda(v-1)$, and counting incidences gives $bk=vr$. Fix a block $B_0$ and let $m_B=\abs{B\cap B_0}$ for each of the other $b-1$ blocks $B$. Each variety of $B_0$ lies in $r-1$ other blocks, and each ordered pair of distinct varieties of $B_0$ lies in $\lambda-1$ other blocks, so
\[
\sum_Bm_B=k(r-1)\quad\text{and}\quad\sum_Bm_B(m_B-1)=k(k-1)(\lambda-1).
\]
Insert these sums into $(b-1)\sum_Bm_B^2\ge\big(\sum_Bm_B\big)^2$, and eliminate $b$ and $\lambda$ with the two identities. A computation shows that the difference of the two sides equals $(r-k)(v-k)(r-\lambda)$. Here $v>k$, and $r>\lambda$ because $r(k-1)=\lambda(v-1)$ with $k<v$; so $r\ge k$, and $b=vr/k\ge v$.
\end{proof}

The proof uses every parameter of the design, and its decisive step is an identity that explains nothing. A relative from geometry, with no parameters at all, is a theorem of de Bruijn and Erd\H{o}s.\index{de Bruijn--Erd\H{o}s theorem}\sidenote{For $n$ points in the real plane, not all on one line, the lines through pairs of them satisfy the hypotheses, so the points determine at least $n$ lines.}

\begin{theorem}[\textcite{deBruijnErdos1948}]\label{thm:dimension:dbe}
Let $X$ be a finite set of $n\ge2$ points and let $\cL$ be a family of $m$ subsets of $X$, called lines, each of which has at least two points, such that every two distinct points lie on exactly one line. If $X$ itself is not a line, then $m\ge n$.
\end{theorem}

\begin{proof}
We give a later double-counting proof rather than the original one. For a point $x$ let $r_x$ be the number of lines through $x$, and for a line $L$ let $k_L=\abs L$. The plan is to average over the pairs $(x,L)$ with $x\notin L$ in two ways.

\emph{Two facts.} If $x\notin L$, then $r_x\ge k_L$, because the lines joining $x$ to the points of $L$ are distinct: a line through $x$ meeting $L$ twice would be $L$. Next, every line misses a point, because $X$ is not a line, and every point $x$ misses a line: otherwise take $y\ne x$ and $z$ off the line $xy$; the line $yz$ contains $x$ and $y$, so it is the line $xy$, which then contains $z$. Thus $m-r_x\ge1$ and $n-k_L\ge1$ for every point $x$ and line $L$.

\emph{The averaging.} Suppose that $m<n$. For $x\notin L$, the first fact and $m<n$ give $nr_x>mk_L$, that is, $n(m-r_x)<m(n-k_L)$. A point $x$ misses exactly $m-r_x$ lines and a line $L$ misses exactly $n-k_L$ points, so summing reciprocals over the pairs with $x\notin L$, first by points and then by lines, gives
\[
1=\sum_{x}\sum_{L\not\ni x}\frac{1}{n(m-r_x)}>\sum_{L}\sum_{x\notin L}\frac{1}{m(n-k_L)}=1,
\]
a contradiction. Hence $m\ge n$.
\end{proof}

The argument is ingenious and fragile: its first fact fails as soon as two points may share two lines. What one would like is a single statement containing both theorems: \emph{a family of distinct subsets of an $n$-set, $n\ge1$, in which every two members meet in exactly $\lambda\ge1$ points has at most $n$ members.} \Cref{sec:dimension:after-independence} proves it in one paragraph.

\subsection{Borsuk's partition problem}\label{sec:dimension:borsuk-before}

For a bounded set $S\subset\R^d$ with at least two points, let $b(S)$ be the least number of parts in a partition of $S$ into sets of diameter smaller than $\diam S$, and let $f(d)$ be the largest $b(S)$ over such $S\subset\R^d$. The $d+1$ vertices of a regular simplex are at mutual distance equal to the diameter, so $f(d)\ge d+1$. In 1933 Borsuk proved that the $d$-dimensional ball also needs $d+1$ parts, using the antipodal theorem, proved in the same paper, that now bears his and Ulam's names (\cref{ch:kneser}). He asked, without guessing the answer, whether $d+1$ parts always suffice \parencite{Borsuk1933}.\index{Borsuk's conjecture}\sidenote{The same paper thus contains the topological lower bound of \cref{ch:kneser} and the partition problem of this chapter.} The positive answer became known as Borsuk's conjecture, and it ``seems to have been believed generally'' \parencite[60]{KahnKalai1993}.

The evidence was geometric. A planar set of diameter $1$ lies in a regular hexagon of width $1$, which three cuts from its center divide into pieces of diameter $\sqrt3/2$. Eggleston settled three-dimensional space in 1955 with a much harder argument, and the conjecture was proved in every dimension for convex bodies with smooth boundary and for centrally symmetric convex bodies. In general dimension the same approach, enclose the set in a convenient body and cut the body, gave only $f(d)\le(\sqrt{3/2}+o(1))^d$; the introduction of \textcite{KahnKalai1993} surveys these results. Erd\H{o}s and Rogers, among others, had suggested that a combinatorial counterexample might exist, but none was found until 1993.

\subsection{Kakeya sets over finite fields}\label{sec:dimension:kakeya-before}

Let $\F_q$ be a field with $q$ elements and $n\ge1$. A \term{line} in $\F_q^n$ is a set $\{a+tb:t\in\F_q\}$ with $b\ne0$; its \term{direction} $b$ is defined up to a nonzero scalar, so there are $(q^n-1)/(q-1)$ directions. A \term{Kakeya set} is a set $K\subseteq\F_q^n$ that contains a line in every direction.\index{Kakeya set!finite field}\sidenote{The name comes from Kakeya's needle problem: how small an area suffices to turn a unit segment around in the plane? Besicovitch showed that it can be made as small as we please.} Wolff asked whether every Kakeya set has at least $c_nq^n$ points, with $c_n>0$ depending only on $n$ \parencite{Wolff1999}. One line gives $\abs K\ge q$, and counting incidences gives more.

\begin{proposition}[the bush bound]\label[proposition]{prop:dimension:bush}
Every Kakeya set $K\subseteq\F_q^n$ satisfies $\abs K^2>q(q^n-1)$; in particular $\abs K>q^{(n+1)/2}/\sqrt2$.
\end{proposition}

\begin{proof}
Choose one line in $K$ for each of the $N=(q^n-1)/(q-1)$ directions. These lines have $qN$ incidences with the points of $K$, so some point $x\in K$ lies on at least $qN/\abs K$ of them. Lines through $x$ with different directions meet only at $x$, so the union of $r$ of them has $1+(q-1)r$ points, all in $K$. Hence $\abs K\ge1+(q-1)qN/\abs K>q(q^n-1)/\abs K$, which is the first claim. The second follows because $q^n-1\ge q^n/2$.
\end{proof}

The bush is a point with many lines through it. Wolff's hairbrush, a line with many lines through it, raised the exponent to $(n+2)/2$ \parencite{Wolff1999}. Arithmetic combinatorics later gave exponents near $4n/7$ in high dimensions and, through the sum-product phenomenon, $5/2+\eps$ in $\F_p^3$; the introduction of \textcite{Dvir2009} surveys these bounds. All of them count incidences, and none reached the exponent $n$.

\subsection{Cap sets}\label{sec:dimension:capset-before}

Three distinct points of $\F_3^n$ lie on a line exactly when they sum to zero: the line $\{x,x+r,x+2r\}$ sums to $3x+3r=0$, and if $x+y+z=0$ and $r=y-x\ne0$, then $z=-x-y=x+2r$. If two of $x,y,z$ are equal and $x+y+z=0$, then all three are equal, because $-2x=x$. So a \term{cap set}, a set $A\subseteq\F_3^n$ with no three distinct points on a line,\sidenote{The word comes from finite geometry, where a cap is a set of points with no three on a line.} is a set in which $x+y+z=0$ has only the solutions $x=y=z$, or equivalently a set with no three-term progression $x,x+r,x+2r$, $r\ne0$ (\cref{fig:dimension:cap}).\index{cap set} The cards of the game SET are the $81$ points of $\F_3^4$, one coordinate for each of four attributes, and a SET is a line.\sidenote{A cap in $\F_3^4$ has at most $20$ points, as Pellegrino proved in 1970, so at most $20$ cards can lie on the table with no SET.}

\begin{figure}[t]
\centering
\begin{tikzpicture}[scale=0.95]
  \foreach \x in {0,1,2} \foreach \y in {0,1,2} {\fill[seamist] (\x,\y) circle (0.13);}
  \draw[sandline,line width=1.1pt] (-0.4,2) -- (2.4,2);
  \draw[sandline,line width=1.1pt] (2,-0.4) -- (2,2.4);
  \draw[sea!60,line width=1.1pt,dashed] (-0.3,0.7) -- (1.3,2.3);
  \draw[sea!60,line width=1.1pt,dashed] (1.7,-0.3) -- (2.3,0.3);
  \foreach \p in {(0,0),(1,0),(0,1),(1,1)} {\fill[seadeep] \p circle (0.13);}
  \foreach \p in {(2,0),(2,1),(0,2),(1,2),(2,2)} {\draw[seadeep,line width=0.6pt] \p circle (0.13);}
  \node[font=\small,muted] at (1,-0.85) {$\F_3^2$};
\end{tikzpicture}
\caption{The plane $\F_3^2$ has nine points and twelve lines of three points each. Drawn are a row, a column, and the broken diagonal $\{(0,1),(1,2),(2,0)\}$, which wraps around. The four filled points form a cap; no five points of $\F_3^2$ do.}\label{fig:dimension:cap}
\end{figure}

Products of caps are caps, so the largest caps grow exponentially with $n$. The best upper bounds before 2016 came from Fourier analysis, through the density increment that Roth had used in 1953 for progressions of integers.\index{density increment}

\begin{aside}[Meshulam's density increment, in outline]
Let $A\subseteq\F_3^n$ be a cap of density $\alpha=\abs A/3^n$, and let $\widehat{1_A}(r)=3^{-n}\sum_{x\in A}\omega^{-r\cdot x}$ with $\omega=e^{2\pi i/3}$. Counting solutions with characters gives
\[
\#\{(x,y,z)\in A^3: x+y+z=0\}=3^{2n}\sum_{r\in\F_3^n}\widehat{1_A}(r)^3 .
\]
For a cap the left side is $\alpha3^n$ and the term $r=0$ on the right is $\alpha^33^{2n}$. So once $3^n\ge2/\alpha^2$, the other terms cancel at least $\alpha^3/2$, and Parseval's identity $\sum_r\abs{\widehat{1_A}(r)}^2=\alpha$ yields some $r\ne0$ with $\abs{\widehat{1_A}(r)}\ge\alpha^2/2$. Then one of the three cosets of the hyperplane $\{x:r\cdot x=0\}$ contains $A$ with relative density at least $\alpha+\alpha^2/4$, and this coset is an affine copy of $\F_3^{n-1}$ in which $A$ is still a cap. Iterate. The density doubles within about $4/\alpha$ steps and never exceeds $1$, so after about $8/\alpha$ steps the condition $3^{n'}\ge2/\alpha'^2$ fails in the current dimension $n'$. Hence $n=O(1/\alpha)$, that is, $\abs A=O(3^n/n)$: each step spends a whole dimension to gain only $\alpha^2/4$ in density, so the bound beats $3^n$ only by a factor of order $n$. See \cref{ch:fourier} for the Fourier dictionary.
\end{aside}

Meshulam proved this bound in 1995 \parencite{Meshulam1995}. Beating it by more than a constant factor took until 2012 and nearly thirty pages, in which Bateman and Katz reached $O(3^n/n^{1+\eps})$ for some $\eps>0$ \parencite{BatemanKatz2012}. The largest caps known had about $2.2174^n$ points \parencite{Edel2004}. Whether every cap has at most $c^n$ points for some $c<3$ was open.

%==========================================================================
\section{The move: count dimensions, not objects}\label{sec:dimension:move}

All five problems ask for the size of a family. The move is to stop counting members and count dimensions: find a vector space in which the members become linearly independent, or in which a small family would carry a polynomial that cannot exist. The engine is one inequality.

\begin{lemma}[rank principle]\label[lemma]{lem:dimension:rank}\index{rank principle}
Let $M$ be an $m\times m$ matrix over a field that factors as $M=BC$, with $B$ of size $m\times D$ and $C$ of size $D\times m$. Then $\rank M\le D$. In particular, if $M$ is nonsingular, then $m\le D$.
\end{lemma}

\begin{proof}
The columns of $M$ are combinations of the $D$ columns of $B$.
\end{proof}

\paragraph{Independence.} The first form of the move turns a hypothesis into a nonsingular matrix. If $v_1,\dots,v_m\in\F^D$ have a nonsingular Gram matrix $(\ip{v_i}{v_j})=VV^{\mathsf T}$, where $V$ has rows $v_i$, then $m\le D$. Likewise, if functions $f_1,\dots,f_m$ lie in a $D$-dimensional space of functions on a set $X$, and $(f_i(x_j))$ is nonsingular for some $x_1,\dots,x_m\in X$, then $m\le D$, because the matrix factors through the coefficients of the $f_i$ in a basis. For set systems the dictionary is one identity: a subset $A$ of an $n$-set becomes its characteristic vector $v_A\in\{0,1\}^n$, and
\begin{equation}\label{eq:dimension:dictionary}
\abs{A\cap B}=\ip{v_A}{v_B}.
\end{equation}
A hypothesis about intersections is a hypothesis about a Gram matrix, whose rank is at most $n$.\index{Gram matrix} The field is ours to choose: $\F_2$ sees parities, the reals see positivity, $\F_p$ sees congruences, and \cref{sec:dimension:fw} uses two fields at once.

\paragraph{Vanishing.} The second form of the move runs the other way. A small set carries a nonzero polynomial of low degree that vanishes on it, for the same dimension count.\index{polynomial method}

\begin{lemma}[interpolation]\label[lemma]{lem:dimension:interpolation}
Let $\F$ be a field and $S\subseteq\F^n$ a finite set. If $\abs S<\binom{n+e}{n}$, then some nonzero polynomial of degree at most $e$ in $n$ variables vanishes at every point of $S$.
\end{lemma}

\begin{proof}
The monomials $x_1^{a_1}\cdots x_n^{a_n}$ with $a_1+\dots+a_n\le e$ form a basis of the polynomials of degree at most $e$, and there are $\binom{n+e}{n}$ of them. The evaluation map from this space to $\F^S$ is linear, and its domain has larger dimension than its target, so its kernel contains a nonzero polynomial.
\end{proof}

The rest of a vanishing argument uses the structure of the set (lines, sums) to make the polynomial vanish at more points than its degree allows; the only fact needed is that a nonzero polynomial of degree $d$ in one variable has at most $d$ roots.

What would one have had to notice? That the answers look like dimensions. The bound $n$ for Oddtown is the dimension of $\F_2^n$, the bound $\binom ms$ of \cref{thm:dimension:fw} is the number of $s$-subsets of an $m$-set, and Dvir's bound $\binom{q+n-1}{n}$ for Kakeya sets is the number of monomials of degree at most $q-1$ in $n$ variables; a binomial coefficient is often a dimension in disguise. Grothendieck's advice to stoop low applies literally. The decisive objects are a column of zeros and ones and a polynomial of degree $q-1$, and nothing about them is deep; what is new is the decision to look at them.

%==========================================================================
\section{After: independence does the counting}\label{sec:dimension:after-independence}

Every dimension bound in this section has the same three steps: encode the sets as vectors, certify that a matrix built from them is nonsingular, and count dimensions.

\subsection{Oddtown and Fisher's inequality}

\begin{theorem}[Oddtown]\label{thm:dimension:oddtown}
Let $\cF$ be a family of subsets of an $n$-element set such that every member of $\cF$ has odd size and every two distinct members meet in an even number of elements. Then $\abs{\cF}\le n$.
\end{theorem}

\begin{proof}
Work over $\F_2$. By \eqref{eq:dimension:dictionary}, $\ip{v_A}{v_B}$ is $\abs{A\cap B}$ modulo $2$, which is $1$ when $A=B$ and $0$ when $A\ne B$. So the Gram matrix of the vectors $v_A$, $A\in\cF$, is the identity, and the rank principle gives $\abs\cF\le\dim\F_2^n=n$.
\end{proof}

Fisher's inequality and the de Bruijn--Erd\H{o}s theorem now take two sentences each.

\begin{proof}[Second proof of \cref{thm:dimension:fisherbibd,thm:dimension:dbe}]
For a design, let $N$ be the incidence matrix, with a row for each variety and a column for each block. By \eqref{eq:dimension:dictionary} and the counts in the proof of \cref{thm:dimension:fisherbibd}, $NN^{\mathsf T}=(r-\lambda)I+\lambda J$ with $r>\lambda$, where $J$ is the all-ones matrix; this matrix is positive definite, so the rank principle gives $v\le b$. For points and lines, $NN^{\mathsf T}=\Delta+J$, where $\Delta$ is diagonal with entries $r_x-1$, and every $r_x\ge2$, since a point on only one line would make that line all of $X$; so $NN^{\mathsf T}$ is positive definite and $n\le m$.
\end{proof}

The first argument is due to \textcite{Bose1949}. Both use only that the Gram matrix is a positive diagonal matrix plus a multiple of $J$, and a little care about the diagonal gives the general statement that \cref{sec:dimension:fisher-before} asked for.

\begin{theorem}[nonuniform Fisher inequality]\label{thm:dimension:fisher}
Let $\lambda\ge1$, and let $\cF$ be a family of distinct subsets of an $n$-element set, $n\ge1$, such that every two distinct members of $\cF$ meet in exactly $\lambda$ elements. Then $\abs\cF\le n$.
\end{theorem}

\begin{proof}
A family with at most one member satisfies the bound because $n\ge1$, so let $\cF=\{A_1,\dots,A_m\}$ with $m\ge2$. By \eqref{eq:dimension:dictionary}, the Gram matrix of the characteristic vectors is $G=\lambda J+\Delta$, where $\Delta$ is diagonal with entries $d_i=\abs{A_i}-\lambda$. Each $d_i\ge0$, because $A_i$ contains its intersection with another member, and at most one $d_i$ is zero: if $\abs{A_i}=\abs{A_j}=\lambda$ with $i\ne j$, then $A_i=A_i\cap A_j=A_j$. For a nonzero $c\in\R^m$,
\[
c^{\mathsf T}Gc=\lambda\Big(\sum_i c_i\Big)^2+\sum_i d_ic_i^2 .
\]
Both terms are nonnegative. If the second vanishes, then $c$ vanishes at every $i$ with $d_i>0$; since $c\ne0$, some $d_{i_0}$ is zero, $c$ is supported on $\{i_0\}$, and the first term is $\lambda c_{i_0}^2>0$ because $\lambda\ge1$. So $G$ is positive definite, and the rank principle gives $m\le n$.
\end{proof}

It contains both earlier theorems. For a design, apply it to the sets of blocks through the varieties, which are distinct because $r>\lambda$; for points and lines, apply it to the sets of lines through the points, which are distinct because $X$ is not a line.

\subsection{Intersections modulo a prime}\label{sec:dimension:fw}

Oddtown prescribes intersections modulo $2$ and Fisher prescribes them exactly. The next theorem prescribes them modulo a prime $p$, and it is here that the move needs two fields at once: the hypotheses live in $\F_p$, while the dimension count is cleanest over $\Q$.\index{Frankl--Wilson theorem}

\begin{theorem}[\textcite{FranklWilson1981}]\label{thm:dimension:fw}
Let $p$ be a prime, let $L$ be a set of $s$ residues modulo $p$, and let $\cF$ be a family of $k$-element subsets of an $m$-element set $X$, where $s\le k$. Suppose that $k$ modulo $p$ does not lie in $L$, and that $\abs{A\cap B}$ modulo $p$ lies in $L$ for every two distinct members $A,B$ of $\cF$. Then $\abs\cF\le\binom ms$.
\end{theorem}

\begin{proof}
The plan is to build an integer matrix indexed by $\cF\times\cF$ that is nonsingular modulo $p$ but has rank at most $\binom ms$ over $\Q$. Choose integers $\ell_1,\dots,\ell_s$ representing the residues in $L$, let $P(t)=\prod_{i=1}^s(t-\ell_i)$, and let $M$ be the matrix with entries $M_{AB}=P(\abs{A\cap B})$.

\emph{Modulo $p$.} For distinct $A,B\in\cF$, the number $\abs{A\cap B}$ is congruent to some $\ell_i$, so $p$ divides $M_{AB}$. On the diagonal, $M_{AA}=P(k)$ is a product of the integers $k-\ell_i$, none of which is divisible by $p$ because $k$ modulo $p$ is not in $L$. So $\det M\equiv P(k)^{\abs\cF}\not\equiv0$ modulo $p$. In particular $\det M\ne0$, and $M$ has rank $\abs\cF$ over $\Q$.

\emph{Over $\Q$.} The polynomials $\binom t0,\dots,\binom ts$ have degrees $0,\dots,s$, so they form a basis of the polynomials of degree at most $s$, and $P(t)=\sum_{j=0}^sa_j\binom tj$ for some rationals $a_j$. The number $\binom{\abs{A\cap B}}{j}$ counts the $j$-subsets $J$ of $A$ with $J\subseteq B$. For such a $J$ and a $k$-set $B$, the $s$-sets $S$ with $J\subseteq S\subseteq B$ number $\binom{k-j}{s-j}$ if $J\subseteq B$ and none otherwise, and $\binom{k-j}{s-j}\ge1$ because $j\le s\le k$. Hence $[J\subseteq B]=\binom{k-j}{s-j}^{-1}\sum_{S}[S\subseteq B]$, the sum running over the $s$-sets $S\supseteq J$. Expand $P$ in the binomial basis, and then replace each indicator $[J\subseteq B]$ by this sum:
\[
M_{AB}=\sum_{j=0}^s a_j\sum_{\substack{J\subseteq A\\ \abs J=j}}[J\subseteq B]=\sum_{\abs S=s}c_{AS}\,[S\subseteq B]
\]
for rationals $c_{AS}$ that do not depend on $B$. This says $M=CW^{\mathsf T}$, where $C=(c_{AS})$ and $W=([S\subseteq B])_{B,S}$ have one column for each of the $\binom ms$ subsets $S$ of size $s$. The rank principle gives $\abs\cF\le\binom ms$.
\end{proof}

\begin{corollary}\label[corollary]{cor:dimension:forbidden}
Let $p$ be a prime and let $\cF$ be a family of $(2p-1)$-element subsets of a $(4p-1)$-element set, no two of which meet in exactly $p-1$ elements. Then $\abs\cF\le\binom{4p-1}{p-1}$.
\end{corollary}

\begin{proof}
Two distinct members meet in between $0$ and $2p-2$ elements, and not in $p-1$, so their intersection size modulo $p$ lies in $L=\{0,1,\dots,p-2\}$, a set of $s=p-1$ residues. The size $k=2p-1$ is congruent to $p-1$, which is not in $L$, and $s\le k$. \Cref{thm:dimension:fw} applies with $m=4p-1$.
\end{proof}

Frankl and Wilson proved the corollary for every prime power in place of $p$ (Exercise~\ref{ex:dimension:primepower}), and used theorems of this kind to show that the chromatic number of $\R^n$, the least number of colors for which no two points at distance $1$ get the same color, grows exponentially in $n$.

\subsection{Borsuk's conjecture is false}\label{sec:dimension:kk}

Kahn and Kalai turned \cref{cor:dimension:forbidden} into a counterexample by realizing ``meet in exactly $p-1$ elements'' as ``realize the diameter''.\sidenote{Their paper, which ended a sixty-year-old conjecture, is three pages long.} The construction has three layers, shown in \cref{fig:dimension:kk}: sets become sign vectors, and sign vectors become points through their tensor squares.

\begin{figure}[t]
\centering
\begin{tikzpicture}[
  box/.style={draw=sea!50,fill=seafoam,rounded corners=1.5pt,inner sep=4pt,font=\small,text width=3.25cm,align=center,minimum height=1.2cm},
  cond/.style={font=\small,text width=3.25cm,align=center},
  arr/.style={-{Stealth[length=2.2mm]},seadeep,line width=0.6pt}]
  \node[box] (s) at (0,0) {a $(2p{-}1)$-subset $A$ of $\{2,\dots,4p\}$};
  \node[box] (x) at (4.05,0) {$x\in\{\pm1\}^{4p}$, \\ $x_i=-1$ iff $i\in\{1\}\cup A$};
  \node[box] (u) at (8.1,0) {$u(x)=(x_ix_j)_{i<j}$ \\ in $\R^{\binom{4p}{2}}$};
  \draw[arr] (s) -- (x); \draw[arr] (x) -- (u);
  \node[cond] (c1) at (0,-1.5) {$\abs{A\cap B}=p-1$};
  \node[cond] (c2) at (4.05,-1.5) {$\ip xy=0$};
  \node[cond] (c3) at (8.1,-1.5) {$\norm{u(x)-u(y)}=\diam$};
  \draw[seadeep,<->,line width=0.5pt] (c1) -- (c2); \draw[seadeep,<->,line width=0.5pt] (c2) -- (c3);
\end{tikzpicture}
\caption{The Kahn--Kalai construction. The forbidden intersection of \cref{cor:dimension:forbidden} becomes orthogonality of sign vectors, and orthogonality becomes the largest distance between tensor squares.}\label{fig:dimension:kk}
\end{figure}

\begin{theorem}[\textcite{KahnKalai1993}]\label{thm:dimension:kahnkalai}
Let $p$ be a prime and $d=\binom{4p}{2}-1$. There is a set of $\binom{4p-1}{2p-1}$ points in $\R^d$ that cannot be partitioned into fewer than $\binom{4p-1}{2p-1}\big/\binom{4p-1}{p-1}$ sets of smaller diameter. For $p=13$ this number exceeds $1561$, so Borsuk's conjecture fails in dimension $1325$.
\end{theorem}

\begin{proof}
The plan is to build the set so that its diameter is attained exactly at the forbidden intersection of \cref{cor:dimension:forbidden}, and then to read off the number of parts. Let $n=4p$, and let $V$ be the set of vectors $x\in\{-1,1\}^n$ with $x_1=-1$ and exactly $2p$ coordinates equal to $-1$. Each $x\in V$ is determined by the set $A_x\subseteq\{2,\dots,n\}$ of its other $2p-1$ coordinates equal to $-1$, so $\abs V=\binom{4p-1}{2p-1}$.

\emph{Inner products.} For $x,y\in V$, the coordinates where $x$ and $y$ differ form the symmetric difference of $\{1\}\cup A_x$ and $\{1\}\cup A_y$, which has $4p-2-2\abs{A_x\cap A_y}$ elements. Since $\ip xy$ is $n$ minus twice the number of such coordinates, $\ip xy=4\big(\abs{A_x\cap A_y}-(p-1)\big)$. So $\ip xy=0$ exactly when $\abs{A_x\cap A_y}=p-1$.

\emph{Tensor squares.} Map $x\in V$ to $u(x)=(x_ix_j)_{1\le i<j\le n}\in\R^{\binom n2}$. Put $w_i=x_iy_i$; then $(x_ix_j-y_iy_j)^2=2-2w_iw_j$, and $2\sum_{i<j}w_iw_j=\ip xy^2-n$ because $\sum_iw_i=\ip xy$ and $w_i^2=1$. Hence
\[
\norm{u(x)-u(y)}^2=n(n-1)-\big(\ip xy^2-n\big)=n^2-\ip xy^2 .
\]
So distances in $u(V)$ are at most $n$, with equality exactly for the pairs with $\ip xy=0$, and such pairs exist (choose $A_x$ and $A_y$ with $p-1$ common elements). The map $u$ is injective on $V$, because $u(x)$ records $x_1x_j$ for every $j$ and $x_1=-1$. Every $u(x)$ lies in the affine hyperplane $\sum_{i<j}z_{ij}=-n/2$, because $\sum_ix_i=0$ gives $2\sum_{i<j}x_ix_j=0-n$; this hyperplane is a copy of $\R^d$.

\emph{Counting parts.} A part of $u(V)$ with diameter smaller than $n$ contains no two points at distance $n$, so the sets $A_x$ of its points are $(2p-1)$-subsets of the $(4p-1)$-set $\{2,\dots,n\}$, no two of which meet in exactly $p-1$ elements. By \cref{cor:dimension:forbidden} a part has at most $\binom{4p-1}{p-1}$ points, which gives the bound on the number of parts. For $p=13$ the quotient is $\binom{51}{25}/\binom{51}{12}=1561.9\ldots$, so at least $1562$ parts are needed, while $d+1=1326$.
\end{proof}

\begin{remark}\label[remark]{rem:dimension:kkgrowth}
The number of parts grows much faster than $d$. Let $H(x)=-x\log_2x-(1-x)\log_2(1-x)$ be the binary entropy. The standard estimates $\binom{4p-1}{2p-1}\ge2^{4p-1}/(4p)$ and $\binom{4p-1}{p-1}\le2^{(4p-1)H(1/4)}$ show that the quotient in \cref{thm:dimension:kahnkalai} is at least $2^{(4p-1)(1-H(1/4))}/(4p)$, which is at least $1.203^{(1-o(1))\sqrt d}$ because $d<8p^2$. Since $f$ is nondecreasing ($\R^d\subset\R^{d+1}$) and consecutive primes have ratio tending to $1$, $f(d)\ge1.2^{\sqrt d}$ for all large $d$, as \textcite{KahnKalai1993} concluded. Monotonicity also spreads the counterexamples: $p=13$ covers $1325\le d\le1560$, $q=16$ (Exercise~\ref{ex:dimension:primepower}) covers $2015\le d\le7501$, and the primes from $17$ on cover every larger $d$, as Kahn and Kalai stated.
\end{remark}

%==========================================================================
\section{After: polynomials that vanish too much}\label{sec:dimension:after-vanishing}

\subsection{Kakeya sets}\label{sec:dimension:dvir}

\begin{lemma}\label[lemma]{lem:dimension:nonvanishing}
Let $\F$ be a field with $q$ elements and let $g$ be a nonzero polynomial in $n$ variables over $\F$ of degree less than $q$. Then $g(x)\ne0$ for some $x\in\F^n$.
\end{lemma}

\begin{proof}
We induct on $n$. A nonzero polynomial in one variable of degree less than $q$ has fewer than $q$ roots, so it does not vanish on all of $\F$. For $n\ge2$, write $g=\sum_{i=0}^jg_i(x_1,\dots,x_{n-1})\,x_n^i$ with $g_j\ne0$. The degree of $g_j$ is at most $\deg g-j<q$, so by induction $g_j(a)\ne0$ for some $a\in\F^{n-1}$. Then $g(a,t)$ is a nonzero polynomial in $t$ of degree $j<q$, and some $t\in\F$ is not a root.
\end{proof}

\begin{theorem}[\textcite{Dvir2009}]\label{thm:dimension:dvir}
Let $\F$ be a field with $q$ elements and let $n\ge1$. Every Kakeya set $K\subseteq\F^n$ has at least $\binom{q+n-1}{n}$ points. In particular $\abs K\ge q^n/n!$.
\end{theorem}

\begin{proof}
Suppose that $\abs K<\binom{q+n-1}{n}=\binom{n+(q-1)}{n}$. By \cref{lem:dimension:interpolation} with $e=q-1$, some nonzero polynomial $g$ of degree $d\le q-1$ vanishes on $K$. The degree $d$ is at least $1$, because $K$ is not empty and a nonzero constant vanishes nowhere. Let $g_d$ be the homogeneous part of $g$ of degree $d$, which is a nonzero polynomial.

Fix a direction $b\ne0$ and a line $\{a+tb:t\in\F\}$ contained in $K$. The polynomial $h(t)=g(a+tb)$ in one variable has degree at most $d\le q-1$ and vanishes at all $q$ elements of $\F$, so it is the zero polynomial. Its coefficient of $t^d$ is $g_d(b)$, because only the monomials of degree $d$ in $g$ contribute to $t^d$. Hence $g_d(b)=0$ for every $b\ne0$, and also $g_d(0)=0$ because $g_d$ is homogeneous of positive degree. So the nonzero polynomial $g_d$, of degree $d<q$, vanishes on all of $\F^n$, which contradicts \cref{lem:dimension:nonvanishing}. The last claim holds because $\binom{q+n-1}{n}=q(q+1)\cdots(q+n-1)/n!$.
\end{proof}

Exercise~\ref{ex:dimension:plane} shows that for odd $q$ the bound is sharp in the plane up to $(q-1)/2$ points; in higher dimensions the right constant is known too (see the Notes).

\subsection{Cap sets and slice rank}\label{sec:dimension:capset}

For the cap set problem the vectors and the polynomial are no longer enough by themselves; the right object is a function of three variables, and the right notion of its size is a rank that treats the three variables symmetrically. For $u\in A$ let $\delta_u$ be the function on $A$ that is $1$ at $u$ and $0$ elsewhere.\index{slice rank}\sidenote{Tao called it simply the rank of a function of several variables; the name slice rank came later.}

\begin{definition}[slice rank]\label[definition]{def:dimension:slicerank}
Let $A$ be a finite set and $\F$ a field. A \term{slice} is a nonzero function $T\colon A\times A\times A\to\F$ of one of the forms $f(x)g(y,z)$, $f(y)g(x,z)$, or $f(z)g(x,y)$. The \term{slice rank} $\dimensionSR(T)$ of a function $T\colon A^3\to\F$ is the least number of slices with sum $T$; the zero function has slice rank $0$.
\end{definition}

For functions of two variables the same definition gives the rank of a matrix, with slices in the role of rank-one matrices. The key lemma says that, as for matrices, a diagonal function has full rank.

\begin{lemma}[\textcite{Tao2016}]\label[lemma]{lem:dimension:diagonal}
Let $A$ be a finite set, $\F$ a field, and $c_u\in\F$ nonzero for every $u\in A$. Then the function $T(x,y,z)=\sum_{u\in A}c_u\delta_u(x)\delta_u(y)\delta_u(z)$ has slice rank $\abs A$.
\end{lemma}

\begin{proof}
The $\abs A$ slices $\delta_u(x)\cdot c_u\delta_u(y)\delta_u(z)$ show that $\dimensionSR(T)\le\abs A$. For the reverse inequality, write $T$ as a sum of $r_1$ slices $f_i(x)g_i(y,z)$, $r_2$ slices $f'_i(y)g'_i(x,z)$, and $r_3$ slices $f''_i(z)g''_i(x,y)$. The plan is to sum against a function of $z$ that kills the slices of the third kind; this turns $T$ into a matrix of large rank and the remaining slices into a matrix of small rank.

\emph{A function with large support.} The functions $h\colon A\to\F$ with $\sum_zh(z)f''_i(z)=0$ for $i=1,\dots,r_3$ form a subspace $H$ of dimension $D\ge\abs A-r_3$. Write a basis of $H$ as the rows of a $D\times\abs A$ matrix; it has rank $D$, so some $D$ of its columns, indexed by a set $S\subseteq A$, are linearly independent. Restriction to $S$ is then a bijection from $H$ to $\F^S$, and the $h\in H$ that equals $1$ on $S$ has at least $D$ nonzero values.

\emph{Contraction.} Multiply the decomposition of $T$ by $h(z)$ and sum over $z$. The left side becomes the matrix $\sum_uc_uh(u)\delta_u(x)\delta_u(y)$ in the variables $x,y$, which is diagonal with at least $D$ nonzero entries, so its rank is at least $D$. On the right, the slices of the third kind contribute zero by the choice of $h$; a slice $f_i(x)g_i(y,z)$ becomes $f_i(x)\sum_zg_i(y,z)h(z)$, a matrix of rank at most $1$, and similarly for the slices of the second kind. Therefore $\abs A-r_3\le D\le r_1+r_2$: every decomposition of $T$ has at least $\abs A$ slices.
\end{proof}

\begin{theorem}[\textcite{EllenbergGijswijt2017}]\label{thm:dimension:capset}
Let $N_n$ be the number of vectors $a\in\{0,1,2\}^n$ with $a_1+\dots+a_n\le2n/3$. Every cap set $A\subseteq\F_3^n$ satisfies $\abs A\le3N_n$. Moreover $N_n\le\Theta^n$, where $\Theta=\min_{0<t\le1}(1+t+t^2)\,t^{-2/3}=2.7551\ldots$
\end{theorem}

\begin{proof}
The plan is to compute the slice rank of the function $T(x,y,z)=[x+y+z=0]$ on $A\times A\times A$, with values in $\F_3$, in two ways.

\emph{The cap condition.} For $x,y,z\in A$, the equation $x+y+z=0$ holds only when $x=y=z$ (\cref{sec:dimension:capset-before}), so $T=\sum_{u\in A}\delta_u(x)\delta_u(y)\delta_u(z)$, and $\dimensionSR(T)=\abs A$ by \cref{lem:dimension:diagonal}.

\emph{A polynomial formula.} For $w\in\F_3$ the number $1-w^2$ is $1$ if $w=0$ and $0$ otherwise, because $w^2=1$ for $w=\pm1$. Hence
\[
T(x,y,z)=\prod_{i=1}^n\big(1-(x_i+y_i+z_i)^2\big).
\]
Expanding the product writes $T$ as a linear combination of monomials $x^ay^bz^c=\prod_ix_i^{a_i}y_i^{b_i}z_i^{c_i}$ with $a_i+b_i+c_i\le2$ for every $i$. Write $\abs a=\sum_ia_i$ for the degree of a multi-index $a$. In each such monomial $a,b,c\in\{0,1,2\}^n$ and $\abs a+\abs b+\abs c\le2n$, so at least one of $\abs a,\abs b,\abs c$ is at most $2n/3$. Assign each monomial to the first of $x,y,z$ whose exponent vector has this property, and collect terms:
\[
T=\sum_{\abs a\le2n/3}x^aF_a(y,z)+\sum_{\abs b\le2n/3}y^bG_b(x,z)+\sum_{\abs c\le2n/3}z^cH_c(x,y)
\]
for some functions $F_a,G_b,H_c$. Each nonzero term is a slice, and each sum has at most $N_n$ terms, so $\dimensionSR(T)\le3N_n$.

\emph{The count.} For $0<t\le1$, every $a\in\{0,1,2\}^n$ with $\abs a\le2n/3$ satisfies $t^{\abs a-2n/3}\ge1$. Summing over all $a\in\{0,1,2\}^n$ instead of these only,
\[
N_n\le\sum_{a\in\{0,1,2\}^n}t^{\abs a-2n/3}=\big((1+t+t^2)\,t^{-2/3}\big)^n,
\]
and minimizing over $t$ gives $N_n\le\Theta^n$.
\end{proof}

The minimum defining $\Theta$ is attained at the positive root $t=(\sqrt{33}-1)/8\approx0.5931$ of $4t^2+t-2=0$, where the derivative of $\log\big((1+t+t^2)t^{-2/3}\big)$ vanishes, and $\Theta=\tfrac38\big(207+33\sqrt{33}\big)^{1/3}=2.755104\ldots$ So every cap set in $\F_3^n$ has fewer than $3\cdot2.7552^n$ points, exponentially fewer than $3^n$.

%==========================================================================
\section{Costs and limits}\label{sec:dimension:costs}

\paragraph{Where did the difficulty go?}
Mostly into the choice of the space and of the field. The engine (rank is at most dimension, and a polynomial of degree $d$ in one variable has at most $d$ roots) was paid for long before these problems were posed, and the one new tool, the slice rank lemma, takes half a page. This is the opposite of the completion move of \cref{ch:lebesgue}, where one hard construction is paid for once and then serves every application. Here each problem needs its own encoding, and finding it is the whole difficulty: once the right space was found, the finite-field Kakeya problem took half a page, but the question waited nine years for it. Is the move then a renaming? No, because the encodings recur: characteristic vectors with a Gram matrix, a polynomial of an intersection size, a vanishing polynomial of degree just below the length of a line, and a monomial count split three ways account for every dimension bound in this chapter, and a new problem pays by choosing from this short list. The price shows elsewhere. A dimension bound sees the size of a family and nothing of its shape; it does not describe the extremal families, it can be far from the truth (Exercise~\ref{ex:dimension:smallcaps}), and it depends on the field, as the prime modulus in \cref{thm:dimension:fw} shows.

\paragraph{What did not dissolve.}
First, the Euclidean Kakeya problem. Dvir's proof uses that a line over $\F_q$ has $q$ points, so a polynomial of degree less than $q$ that vanishes on it vanishes identically along it. In $\R^n$ a Kakeya set, a compact set containing a unit segment in every direction, can have measure zero, and the conjecture is that it has Hausdorff dimension $n$; no finite point set carries the problem. Wang and Zahl announced a proof for $n=3$ in 2025, with multiscale geometric arguments rather than a dimension count \parencite{WangZahl2025}, and $n\ge4$ is open. Second, the cap set exponent. Caps with $2.2202^n$ points exist for large $n$ \parencite{RomeraParedesEtAl2024}, and the best upper bound is still $\Theta^n=2.7551\ldots^n$, up to polynomial factors. Slice rank alone cannot close the gap: the argument applies verbatim to \term{tri-colored sum-free sets},\index{tri-colored sum-free set} families of triples $(a_i,b_i,c_i)$ with $a_i+b_j+c_k=0$ exactly when $i=j=k$, and for those the base $\Theta$ is the truth \parencite{KleinbergSawinSpeyer2018}, with a final step supplied by Pebody and, independently, by Norin. An improvement for caps must use that a cap is one set rather than three. Third, Borsuk's problem in low dimensions: in the refereed literature it is open in every dimension from $4$ to $63$. Fourth, the integers: in $\{1,\dots,N\}$ there is no vector space to count in, and Behrend's 1946 construction of progression-free sets of size $N^{1-o(1)}$ rules out bounds of the form $N^c$ with $c<1$. Finally, a dissolution can also dissolve a hope: slice rank shows that the group-theoretic approach to fast matrix multiplication cannot reach exponent $2$ with abelian groups of bounded exponent \parencite{BlasiakEtAl2017}.

\paragraph{Other seas.}
Fourier analysis (\cref{ch:fourier}) is also linear algebra over a group, but it diagonalizes and measures where the dimension argument only counts; density increments remain the main tool for progressions in the integers. Topology (\cref{ch:kneser}) supplies other certificates: Borsuk's lower bound for the ball comes from the antipodal theorem. Within the polynomial method, the Combinatorial Nullstellensatz (Exercise~\ref{ex:dimension:cd}), polynomials vanishing to high order, and the polynomial partitioning of Guth and Katz carried the vanishing form of the move into additive number theory and incidence geometry \parencite{Guth2016}.

%==========================================================================
\section{The move, distilled}\label{sec:dimension:distilled}

\begin{move}{Linearization}
\movepart{In one sentence} To bound the size of a family, map its members to vectors or functions in a space of known dimension, so that the defining property makes them linearly independent, or makes a small family carry a nonzero low-degree polynomial that the property forces to vanish where it cannot.
\movepart{Triggers} Intersections, inner products, or distances that take few values; parity or congruence conditions; a set that contains a line in every direction or avoids the solutions of a linear equation; an answer that looks like a dimension, such as $n$, $\binom ms$, or a number of monomials.
\movepart{Recipe} (1) Encode each object as a vector (characteristic vector, sign vector, tensor square) or as a function (a polynomial of an inner product). (2) Choose the field in which the hypothesis makes the evaluation matrix diagonal, triangular, or positive definite. (3) Bound the dimension of a space containing the encodings: count monomials, reduce on the cube or on a slice. (4) Apply the rank principle. In the vanishing form, interpolate a low-degree polynomial on the set and let the structure of the set make it vanish too often; for equations in three variables, use slice rank.
\movepart{Anti-pattern} An encoding into a space whose dimension is the trivial bound: the indicator functions of all subsets of an $n$-set span a space of dimension $2^n$, and ``at most $2^n$'' was free. A Gram matrix whose nonsingularity is asserted rather than certified. Congruences modulo a number such as $6$, where no single prime sees the condition and the argument breaks down.
\end{move}

\dissolutiontest{Fisher's inequality and the de Bruijn--Erd\H{o}s theorem go from a page of counting to two sentences each; the Kakeya exponent jumps from about $4n/7$ to $n$ in half a page; the cap set bound goes from an iterated density increment to a one-page rank count, with an exponential gain.}{Fisher's inequality without uniformity, for blocks of different sizes and every $\lambda\ge1$; intersections prescribed only modulo a prime.}{Oddtown, Fisher, Frankl--Wilson, Kahn--Kalai, Dvir, and Ellenberg--Gijswijt are one inequality, rank at most dimension, applied to a Gram matrix, an evaluation matrix, or a slice decomposition.}{Borsuk's conjecture refuted; the finite-field Kakeya conjecture proved; cap sets exponentially small; the chromatic number of $\R^n$ exponential in $n$.}

Linear independence is the notion that \cref{ch:matroid} axiomatizes, and the rank principle is its most quantitative consequence. The VC dimension of \cref{ch:vc} meets this chapter in Exercise~\ref{ex:dimension:vc}: a $D$-dimensional space of functions cannot shatter $D+1$ points, by the same rank count. Borsuk's 1933 paper contains both the antipodal theorem that \cref{ch:kneser} uses for lower bounds and the partition problem that linear algebra refuted. The finite-field Kakeya problem also arose in theoretical computer science, in the analysis of randomness mergers, relatives of the extractors of \cref{ch:pseudorandom}.

%==========================================================================
\begin{exercises}
\item \consolidation\hard Let $\cF$ be a family of subsets of an $n$-element set in which every member, and every intersection of two members, has even size. Prove that $\abs\cF\le2^{\floor{n/2}}$, so the Eventown construction of \cref{sec:dimension:oddtown-before} is optimal. \hint{The characteristic vectors span a subspace $W$ of $\F_2^n$ with $\ip w{w'}=0$ for all $w,w'\in W$; compare $\dim W$ with $\dim W^\perp=n-\dim W$.}\label{ex:dimension:eventown}

\item \consolidation Let $\F$ be a field, $S\subseteq\F$ a finite set, and $g$ a nonzero polynomial of degree $d$ in $n$ variables over $\F$. Prove that $g$ vanishes at no more than $d\abs S^{n-1}$ points of $S^n$ (the Schwartz--Zippel lemma), and deduce \cref{lem:dimension:nonvanishing}.\label{ex:dimension:sz}

\item \consolidation\hard (a) Let $L$ be a set of $s$ integers in $\{0,1,\dots,k-1\}$, and let $\cF$ be a family of $k$-subsets of an $m$-set such that $\abs{A\cap B}\in L$ for every two distinct members. Prove that $\abs\cF\le\binom ms$ (the Ray-Chaudhuri--Wilson theorem). (b) Drop the condition that the members have equal size, and let $L$ be a set of $s$ integers. Prove that $\abs\cF\le\sum_{i=0}^s\binom mi$. \hint{For (a), run the proof of \cref{thm:dimension:fw} over $\Q$ alone. For (b), list the members by increasing size, let $f_A(x)=\prod_{\ell\in L,\,\ell<\abs A}(\ip{v_A}x-\ell)$, check that $(f_A(v_B))$ is triangular with nonzero diagonal, and replace each $x_i^2$ by $x_i$, which changes no value on $\{0,1\}^m$.}\label{ex:dimension:rw}

\item \consolidation\hard (a) Check that the proof of \cref{thm:dimension:fw} uses only this: $M_{AB}=P(\abs{A\cap B})$ for a polynomial $P$ of degree $s\le k$ with rational coefficients and integer values at integers, and $M$ is diagonal modulo $p$ with diagonal entries prime to $p$. (b) Let $q$ be a power of the prime $p$ and $P(t)=\binom{t-q}{q-1}$. Use Lucas's theorem to show that $p$ divides $P(t)$ for every integer $0\le t\le2q-2$ other than $t=q-1$, while $P(2q-1)=1$, and deduce \cref{cor:dimension:forbidden} with $q$ in place of $p$. (c) With $q=16$, show that Borsuk's conjecture fails in dimension $2015$. \hint{For $0\le t\le q-2$, write $P(t)=\pm\binom{2q-2-t}{q-1}$ and compare base-$p$ digits.}\label{ex:dimension:primepower}

\item \consolidation Let $q$ be odd. For $a\in\F_q$ let $\ell_a=\{(x,2ax-a^2):x\in\F_q\}$, and let $K$ be the union of these $q$ lines and the line $\{(0,y):y\in\F_q\}$. Show that $K$ is a Kakeya set with $\abs K=q(q+1)/2+(q-1)/2$, so \cref{thm:dimension:dvir} is sharp in the plane up to $(q-1)/2$ points. \hint{A point $(x,y)$ lies on some $\ell_a$ exactly when $x^2-y$ is a square in $\F_q$, zero included.}\label{ex:dimension:plane}

\item \consolidation\hard (a) Let $p$ be an odd prime and $A\subseteq\F_p^n$ a set with no three-term progression $x,x+r,x+2r$, $r\ne0$. Prove that $\abs A$ is at most $3$ times the number of $a\in\{0,1,\dots,p-1\}^n$ with $\sum_ia_i\le(p-1)n/3$. (b) For $n\le4$, compare $3N_n$ in \cref{thm:dimension:capset} with the largest caps in $\F_3^n$, which have $2$, $4$, $9$, and $20$ points. \hint{For (a), use $T(x,y,z)=[x-2y+z=0]$ and $[w=0]=1-w^{p-1}$ on $\F_p$.}\label{ex:dimension:smallcaps}

\item \transfer\hard A set $S\subset\R^d$ is a \term{two-distance set} if the distances between distinct points of $S$ take only two values $\beta$ and $\gamma$. Prove that $\abs S\le(d+1)(d+4)/2$, a theorem of Larman, Rogers, and Seidel. \hint{For $s\in S$ let $f_s(x)=(\norm{x-s}^2-\beta^2)(\norm{x-s}^2-\gamma^2)$; every $f_s$ lies in the span of $\norm x^4$, the functions $\norm x^2x_i$, the products $x_ix_j$ with $i\le j$, the coordinates $x_i$, and $1$.} The midpoints of the edges of a regular simplex give $\binom{d+1}{2}$ points.\label{ex:dimension:twodistance}

\item \transfer Suppose that the edges of the complete graph $K_n$ are partitioned into the edge sets of $m$ complete bipartite graphs. Prove that $m\ge n-1$, and that $n-1$ can occur (the Graham--Pollak theorem). \hint{This is Tverberg's proof. If the $k$th graph has sides $X_k$ and $Y_k$, then $\sum_{i<j}x_ix_j=\sum_k\big(\sum_{i\in X_k}x_i\big)\big(\sum_{j\in Y_k}x_j\big)$. If $m<n-1$, then the equations $\sum_ix_i=0$ and $\sum_{i\in X_k}x_i=0$ for all $k$ have a nonzero real solution; compute $(\sum_ix_i)^2$ there.}\label{ex:dimension:gp}

\item \transfer Let $\cG$ be a $D$-dimensional vector space of real functions on a set $X$. Prove that the sets $\{x:g(x)\ge0\}$, $g\in\cG$, cannot cut out every subset of a set of $D+1$ points of $X$; in the language of \cref{ch:vc}, their VC dimension is at most $D$. Deduce that half-spaces in $\R^k$ have VC dimension at most $k+1$, and that the sets $\{p\ge0\}$ with $p$ a polynomial of degree at most $e$ in $k$ variables have VC dimension at most $\binom{k+e}{k}$. \hint{If $x_1,\dots,x_{D+1}$ were shattered, the vectors $(g(x_1),\dots,g(x_{D+1}))$ would lie in a hyperplane $\sum_i\lambda_iy_i=0$ with some $\lambda_i<0$; take $g$ with $g(x_i)\ge0$ exactly when $\lambda_i\ge0$.} For homogeneous half-spaces and points in general position, Cover's function-counting theorem of 1965 gives the exact number of subsets cut out.\label{ex:dimension:vc}

\item \transfer\hard The Combinatorial Nullstellensatz \parencite{Alon1999} states: if $f$ is a polynomial over a field of degree $t_1+\dots+t_n$ in which the monomial $x_1^{t_1}\cdots x_n^{t_n}$ has a nonzero coefficient, and $S_1,\dots,S_n$ are sets in the field with $\abs{S_i}>t_i$, then $f(s)\ne0$ for some $s\in S_1\times\dots\times S_n$. Use it to prove the Cauchy--Davenport theorem: nonempty sets $A,B\subseteq\Z/p\Z$, with $p$ prime, satisfy $\abs{A+B}\ge\min(p,\abs A+\abs B-1)$. \hint{If $\abs A+\abs B>p$, show directly that $A+B=\Z/p\Z$. Otherwise suppose $A+B\subseteq C$ with $\abs C=\abs A+\abs B-2$, and apply the theorem to $\prod_{c\in C}(x+y-c)$.}\label{ex:dimension:cd}

\item \transfer\harder Let $\cL$ be a set of $L$ lines in $\R^3$. A \term{joint} is a point that lies on three lines of $\cL$ with linearly independent directions. Prove that there are at most $CL^{3/2}$ joints, for an absolute constant $C$ (the joints theorem of Guth and Katz; see \textcite{Guth2016}). \hint{(i) If there are $J$ joints, some nonzero polynomial of degree at most $(6J)^{1/3}+1$ vanishes at all of them; take one, $P$, of least degree. (ii) If every line of $\cL$ contains more than $\deg P$ joints, then $P$ vanishes on every line, its gradient vanishes at every joint, and some partial derivative of $P$ contradicts the minimality of $\deg P$. (iii) So some line contains at most $\deg P$ joints; delete it and induct.}\label{ex:dimension:joints}

\item \research\harder Let $r(n)$ be the largest size of a cap set in $\F_3^n$. (a) Show that $r(m+n)\ge r(m)r(n)$, and deduce that $c=\lim_nr(n)^{1/n}$ exists and equals $\sup_nr(n)^{1/n}$. (b) The bounds $2.2202\le c\le2.7551\ldots$ are known (\cref{thm:dimension:capset}, \textcite{RomeraParedesEtAl2024}); the value of $c$ is open. Since slice rank alone cannot do better for tri-colored sum-free sets (\cref{sec:dimension:costs}), find a property that distinguishes a cap $A$, viewed as the triples $(a,a,a)$, from a general tri-colored sum-free set, and that a rank argument can detect.\label{ex:dimension:capexponent}

\item \research\harder (a) Let $y_1,\dots,y_N$ be real polynomials in $W<N$ variables. Show that some nonzero polynomial $Q$ in $N$ variables satisfies $Q(y_1,\dots,y_N)=0$ identically, and deduce that a model $f(\theta,x)$ that is polynomial in its $W$ parameters $\theta$ can fit exactly only a Lebesgue-null set of target vectors on $N>W$ inputs. \hint{Compare the dimension of the polynomials of degree at most $r$ in $N$ variables with that of the polynomials of degree at most $re$ in $W$ variables, $e=\max_i\deg y_i$, as $r\to\infty$.} (b) For classification only signs matter. Warren's theorem bounds the number of sign patterns of $N$ polynomials in $W$ variables, so fitting all $2^N$ labellings is also limited by a count, but a weaker one than in (a), and deep networks are known to fit labellings with far fewer parameters than data points. Find dimension-type invariants that predict how many parameters a network of given depth needs for data of given geometry, and decide when they are tight; beyond special cases such as Exercise~\ref{ex:dimension:vc}, this is open.\label{ex:dimension:memorization}
\end{exercises}

\notesandrefs
The linear algebra method in combinatorics is the subject of \textcite{BabaiFrankl1992}, which opens with Oddtown; \textcite{Matousek2010} gives short chapters on Fisher's inequality, the Graham--Pollak theorem, two-distance sets, and the Kahn--Kalai construction. The counting proof of \cref{thm:dimension:fisherbibd} is Fisher's \parencite{Fisher1940}, the matrix proof is due to \textcite{Bose1949}, and the nonuniform version was noted by Majumdar in 1953 and rediscovered by Isbell in 1959. The original proof of \textcite{deBruijnErdos1948} is a different counting argument from ours. Ray-Chaudhuri and Wilson proved the bound of Exercise~\ref{ex:dimension:rw}(a) in 1975, and \textcite{FranklWilson1981} proved the modular theorem, its prime-power form, and, with parameters other than those of \cref{cor:dimension:forbidden}, the bound $(1+o(1))1.2^n$ for the chromatic number of $\R^n$.

Borsuk's question shares its paper with the antipodal theorem \parencite{Borsuk1933}. \textcite{KahnKalai1993} phrased their construction with the edge sets of complete bipartite graphs, which are the tensor squares of \cref{fig:dimension:kk} in disguise. The count printed in their paper is smaller than ours by a factor of $2$, too small to give the dimension $1325$; the count of \cref{thm:dimension:kahnkalai} gives it. The dimension of a counterexample was lowered to $946$ by Nilli, $561$ by Raigorodskii, $298$ by Hinrichs and Richter, $65$ by Bondarenko, and $64$ by \textcite{JenrichBrouwer2014}, the last two with two-distance sets. Preprints posted in 2026 claim counterexamples in lower dimensions and await independent checking.

Wolff proposed the finite-field Kakeya problem as a model for the Euclidean one \parencite{Wolff1999}. The first version of Dvir's paper, posted in March 2008, proved $\abs K\ge C_nq^{n-1}$; Alon and Tao independently observed how the argument gives the bound of \cref{thm:dimension:dvir}, and the published paper includes their improvement \parencite{Dvir2009}. Polynomials vanishing to high order on $K$ (the method of multiplicities) later improved the constant $1/n!$ to about $2^{-n}$, and \textcite{BukhChao2021} found the sharp density, close to $2^{1-n}$. \textcite{Guth2016} follows the polynomial method from Dvir's proof to the joints problem and beyond.

Croot, Lev, and Pach posted their bound for $(\Z/4\Z)^n$ in May 2016 \parencite{CrootLevPach2017}; within two weeks Ellenberg and Gijswijt, independently of each other, adapted it to $\F_q^n$ and wrote one joint paper \parencite{EllenbergGijswijt2017}. The symmetric formulation with slice rank that we followed appeared on Tao's blog on 18 May 2016 \parencite{Tao2016}. Edel's caps \parencite{Edel2004} held the record for nearly twenty years, until computer-assisted searches raised the base to $2.2202$ \parencite{RomeraParedesEtAl2024}.

\printchapterbib
\end{refsection}
